\documentclass[../main.tex]{subfiles}

\begin{document}

\section{Chapter 3: Pre-Attentive Processing, Gestalt Principles, and Magnitude Estimation}
\label{sec:chapter3-preattentive-gestalt}

\subsection{Techniques / Principles Applied}

\begin{itemize}[leftmargin=*]
\item Pre-attentive processing through contrast, saturation, and enclosure.
\item Gestalt grouping, especially proximity, similarity, and enclosure.
\item Magnitude estimation that favors common baselines and linear length for
important comparisons.
\end{itemize}

\subsection{How Applied in the Dashboard}

Pre-attentive pop-out is used to direct attention before detailed reading. On
Page 1, Figure~\ref{fig:socioeconomic-matrix} uses darker blue cells to make
high-density income--education intersections stand out quickly. On Page 2, the
mobile ownership chart differentiates smartphone and basic phone rates with
strong luminance contrast, so the penetration gap is visible without a full
label-by-label scan.

\dashboardfigure{0.86}{fig-page3-kpi-cards.png}
{The Account Ownership \& Access Landscape KPI cards.}
{fig:account-kpi-cards}

Gestalt grouping is especially visible on Page 3. Figure~\ref{fig:account-kpi-cards}
clusters the four KPI cards closely together, so proximity signals that they
belong to the same summary layer. Enclosure appears in the filter strip, where
the controls are visually separated from the charts they drive.

\dashboardfigure{0.92}{fig-page3-filter-bar.png}
{Filter controls on the Account Ownership \& Access Landscape page.}
{fig:page3-filter-bar}

Magnitude estimation is handled conservatively in the comparison-heavy charts.
The Mobile Ownership Profile and the Exclusion Friction Log rely on bar length
and aligned baselines because those channels are more accurate for comparing
rates than area- or angle-based alternatives.

\subsection{Notes / Adjustments}

The Account Composition Flow in Figure~\ref{fig:account-composition-flow} is
less precise than a standard bar chart for direct magnitude comparison. The
trade-off is accepted because the page needs to emphasize relationships between
source groups and account states, not only isolated totals.

\end{document}
